% Options for packages loaded elsewhere
\PassOptionsToPackage{unicode}{hyperref}
\PassOptionsToPackage{hyphens}{url}
\documentclass[
]{article}
\usepackage{xcolor}
\usepackage{amsmath,amssymb}
\setcounter{secnumdepth}{-\maxdimen} % remove section numbering
\usepackage{iftex}
\ifPDFTeX
  \usepackage[T1]{fontenc}
  \usepackage[utf8]{inputenc}
  \usepackage{textcomp} % provide euro and other symbols
\else % if luatex or xetex
  \usepackage{unicode-math} % this also loads fontspec
  \defaultfontfeatures{Scale=MatchLowercase}
  \defaultfontfeatures[\rmfamily]{Ligatures=TeX,Scale=1}
\fi
\usepackage{lmodern}
\ifPDFTeX\else
  % xetex/luatex font selection
\fi
% Use upquote if available, for straight quotes in verbatim environments
\IfFileExists{upquote.sty}{\usepackage{upquote}}{}
\IfFileExists{microtype.sty}{% use microtype if available
  \usepackage[]{microtype}
  \UseMicrotypeSet[protrusion]{basicmath} % disable protrusion for tt fonts
}{}
\makeatletter
\@ifundefined{KOMAClassName}{% if non-KOMA class
  \IfFileExists{parskip.sty}{%
    \usepackage{parskip}
  }{% else
    \setlength{\parindent}{0pt}
    \setlength{\parskip}{6pt plus 2pt minus 1pt}}
}{% if KOMA class
  \KOMAoptions{parskip=half}}
\makeatother
\setlength{\emergencystretch}{3em} % prevent overfull lines
\providecommand{\tightlist}{%
  \setlength{\itemsep}{0pt}\setlength{\parskip}{0pt}}

\usepackage{newunicodechar}
\newunicodechar{❌}{-}

% --------------------------------------------------
% Margini pagina (del testo)
% --------------------------------------------------
\usepackage[
  a4paper,
  left=1.5cm,
  right=1.5cm,
  top=1.5cm,
  bottom=1.5cm
]{geometry}

% --------------------------------------------------
% Immagini
% --------------------------------------------------
\usepackage{graphicx}

% --------------------------------------------------
% Font per XeLaTeX
% --------------------------------------------------
\usepackage{fontspec}
\setmainfont{Times New Roman}

% --------------------------------------------------
% Paragrafi
% --------------------------------------------------
\setlength{\parindent}{0pt}
\setlength{\parskip}{6pt}

% --------------------------------------------------
% Link
% --------------------------------------------------
\usepackage{hyperref}

% --------------------------------------------------
% Sfondo pagina a tutto schermo (senza margini)
% --------------------------------------------------
\usepackage{background}

\backgroundsetup{
  scale=1,
  angle=0,
  opacity=1,
  position=current page.center, % Centra l'immagine sulla pagina fisica
  nodeanchor=center,            % Forza l'ancoraggio al centro esatto
  pages=all,
  contents={
    \includegraphics[
      width=\paperwidth,        % Larghezza totale del foglio
      height=\paperheight,      % Altezza totale del foglio
      keepaspectratio=false     % Forza la copertura totale
    ]{./005_authoring_gen/imgs/sfondocopyright.png}
  }
}
\usepackage{bookmark}
\IfFileExists{xurl.sty}{\usepackage{xurl}}{} % add URL line breaks if available
\urlstyle{same}
\hypersetup{
  hidelinks,
  pdfcreator={LaTeX via pandoc}}

\author{}
\date{}

\begin{document}

\begin{center}\rule{0.5\linewidth}{0.5pt}\end{center}

\section{ESERCIZI CON SOLUZIONI -- LIVELLO 1
(Java)}\label{esercizi-con-soluzioni-livello-1-java}

\begin{center}\rule{0.5\linewidth}{0.5pt}\end{center}

\section{1. Introduzione a Java}\label{introduzione-a-java}

\subsection{Esercizio 1}\label{esercizio-1}

Scrivere un programma che stampi una frase.

\textbf{Soluzione}

\begin{verbatim}
public class Esercizio1 {
    public static void main(String[] args) {
        System.out.println("Benvenuto in Java!");
    }
}
\end{verbatim}

\begin{center}\rule{0.5\linewidth}{0.5pt}\end{center}

\subsection{Esercizio 2}\label{esercizio-2}

Stampare due frasi su due righe diverse.

\textbf{Soluzione}

\begin{verbatim}
public class Esercizio2 {
    public static void main(String[] args) {
        System.out.println("Prima riga");
        System.out.println("Seconda riga");
    }
}
\end{verbatim}

\begin{center}\rule{0.5\linewidth}{0.5pt}\end{center}

\subsection{Esercizio 3}\label{esercizio-3}

Compilare manualmente con javac ed eseguire con java.

\textbf{Soluzione attesa (terminal output)}

\begin{itemize}
\tightlist
\item
  Compilazione: \texttt{javac\ Esercizio3.java}
\item
  Esecuzione: \texttt{java\ Esercizio3}
\end{itemize}

\begin{center}\rule{0.5\linewidth}{0.5pt}\end{center}

\subsection{Esercizio 4}\label{esercizio-4}

Programma che stampa tre istruzioni tipo manuale.

\textbf{Soluzione}

\begin{verbatim}
public class Esercizio4 {
    public static void main(String[] args) {
        System.out.println("1. Accendere il dispositivo");
        System.out.println("2. Attendere il caricamento");
        System.out.println("3. Selezionare una funzione");
    }
}
\end{verbatim}

\begin{center}\rule{0.5\linewidth}{0.5pt}\end{center}

\subsection{Esercizio 5}\label{esercizio-5}

Stampare definizioni di JVM, JRE e JDK.

\textbf{Soluzione}

\begin{verbatim}
public class Esercizio5 {
    public static void main(String[] args) {
        System.out.println("JVM: esegue il bytecode");
        System.out.println("JRE: ambiente per eseguire Java");
        System.out.println("JDK: strumenti di sviluppo, compilatore incluso");
    }
}
\end{verbatim}

\begin{center}\rule{0.5\linewidth}{0.5pt}\end{center}

\section{2. Ambiente di sviluppo}\label{ambiente-di-sviluppo}

\subsection{Esercizio 1}\label{esercizio-1-1}

Verificare le versioni.

\textbf{Soluzione attesa}

\begin{itemize}
\tightlist
\item
  \texttt{java\ -version}
\item
  \texttt{javac\ -version}
\end{itemize}

\begin{center}\rule{0.5\linewidth}{0.5pt}\end{center}

\subsection{Esercizio 4}\label{esercizio-4-1}

Compilare generando file .class in una cartella bin.

\textbf{Comandi}

\begin{itemize}
\tightlist
\item
  \texttt{javac\ -d\ bin\ src/Esempio.java}
\item
  \texttt{java\ -cp\ bin\ Esempio}
\end{itemize}

\begin{center}\rule{0.5\linewidth}{0.5pt}\end{center}

\section{3. Sintassi di base}\label{sintassi-di-base}

\subsection{Esercizio 1}\label{esercizio-1-2}

Classe con main che stampa un messaggio.

\textbf{Soluzione}

\begin{verbatim}
public class TestSintassi {
    public static void main(String[] args) {
        System.out.println("Messaggio di prova");
    }
}
\end{verbatim}

\begin{center}\rule{0.5\linewidth}{0.5pt}\end{center}

\subsection{Esercizio 4}\label{esercizio-4-2}

Due metodi statici richiamati dal main.

\textbf{Soluzione}

\begin{verbatim}
public class MetodiStatici {
    public static void saluta() {
        System.out.println("Ciao!");
    }

    public static void addio() {
        System.out.println("Arrivederci!");
    }

    public static void main(String[] args) {
        saluta();
        addio();
    }
}
\end{verbatim}

\begin{center}\rule{0.5\linewidth}{0.5pt}\end{center}

\section{4. Tipi di dato e variabili}\label{tipi-di-dato-e-variabili}

\subsection{Esercizio 1}\label{esercizio-1-3}

Dichiarare variabili di tutti i tipi primitivi.

\textbf{Soluzione}

\begin{verbatim}
public class TipiPrimitivi {
    public static void main(String[] args) {
        byte a = 10;
        short b = 200;
        int c = 3000;
        long d = 40000L;
        float e = 3.14f;
        double f = 2.71828;
        char g = 'A';
        boolean h = true;

        System.out.println(a + " " + b + " " + c + " " + d);
        System.out.println(e + " " + f + " " + g + " " + h);
    }
}
\end{verbatim}

\begin{center}\rule{0.5\linewidth}{0.5pt}\end{center}

\subsection{Esercizio 4}\label{esercizio-4-3}

Convertire double in int.

\textbf{Soluzione}

\begin{verbatim}
public class Cast {
    public static void main(String[] args) {
        double x = 12.9;
        int y = (int) x;
        System.out.println(y); // stampa 12
    }
}
\end{verbatim}

\begin{center}\rule{0.5\linewidth}{0.5pt}\end{center}

\section{5. Operatori}\label{operatori}

\subsection{Esercizio 1}\label{esercizio-1-4}

Operazioni aritmetiche.

\textbf{Soluzione}

\begin{verbatim}
public class Operatori1 {
    public static void main(String[] args) {
        int a = 12;
        int b = 5;

        System.out.println("Somma: " + (a + b));
        System.out.println("Diff: " + (a - b));
        System.out.println("Prod: " + (a * b));
        System.out.println("Resto: " + (a % b));
    }
}
\end{verbatim}

\begin{center}\rule{0.5\linewidth}{0.5pt}\end{center}

\subsection{Esercizio 6}\label{esercizio-6}

Espressione complessa.

\textbf{Soluzione}

\begin{verbatim}
int risultato = (4 + 5 * 2) - (10 / 2);
\end{verbatim}

\begin{center}\rule{0.5\linewidth}{0.5pt}\end{center}

\section{6. Strutture di selezione}\label{strutture-di-selezione}

\subsection{Esercizio 1}\label{esercizio-1-5}

Verificare pari/dispari.

\textbf{Soluzione}

\begin{verbatim}
public class PariDispari {
    public static void main(String[] args) {
        int n = 7;
        if (n % 2 == 0) {
            System.out.println("Pari");
        } else {
            System.out.println("Dispari");
        }
    }
}
\end{verbatim}

\begin{center}\rule{0.5\linewidth}{0.5pt}\end{center}

\subsection{Esercizio 4}\label{esercizio-4-4}

Calcolare se una persona è maggiorenne.

\textbf{Soluzione}

\begin{verbatim}
public class Maggiorenne {
    public static void main(String[] args) {
        int anno = 2006;
        int eta = 2025 - anno;

        if (eta >= 18) {
            System.out.println("Maggiorenne");
        } else {
            System.out.println("Minorenne");
        }
    }
}
\end{verbatim}

\begin{center}\rule{0.5\linewidth}{0.5pt}\end{center}

\subsection{Esercizio 6}\label{esercizio-6-1}

Simulare accesso a un servizio.

\textbf{Soluzione}

\begin{verbatim}
public class Accesso {
    public static void main(String[] args) {
        String ruolo = "docente";

        if (ruolo.equals("studente")) {
            System.out.println("Accesso limitato");
        } else if (ruolo.equals("docente")) {
            System.out.println("Accesso a sezioni didattiche");
        } else if (ruolo.equals("amministratore")) {
            System.out.println("Accesso completo");
        } else {
            System.out.println("Ruolo non riconosciuto");
        }
    }
}
\end{verbatim}

\begin{center}\rule{0.5\linewidth}{0.5pt}\end{center}

\section{7. Strutture iterative}\label{strutture-iterative}

\subsection{Esercizio 1}\label{esercizio-1-6}

Stampare 1--10.

\textbf{Soluzione}

\begin{verbatim}
for (int i = 1; i <= 10; i++) {
    System.out.println(i);
}
\end{verbatim}

\begin{center}\rule{0.5\linewidth}{0.5pt}\end{center}

\subsection{Esercizio 4}\label{esercizio-4-5}

Somma dei primi n numeri.

\textbf{Soluzione}

\begin{verbatim}
int n = 5;
int somma = 0;
for (int i = 1; i <= n; i++) {
    somma += i;
}
System.out.println(somma);
\end{verbatim}

\begin{center}\rule{0.5\linewidth}{0.5pt}\end{center}

\subsection{Esercizio 6}\label{esercizio-6-2}

Triangolo di asterischi.

\textbf{Soluzione}

\begin{verbatim}
for (int r = 1; r <= 5; r++) {
    for (int c = 1; c <= r; c++) {
        System.out.print("*");
    }
    System.out.println();
}
\end{verbatim}

\begin{center}\rule{0.5\linewidth}{0.5pt}\end{center}

\section{8. Array}\label{array}

\subsection{Esercizio 1}\label{esercizio-1-7}

Creare array e stamparlo.

\textbf{Soluzione}

\begin{verbatim}
int[] a = {3, 7, 2, 9, 5};
for (int i = 0; i < a.length; i++) {
    System.out.println(a[i]);
}
\end{verbatim}

\begin{center}\rule{0.5\linewidth}{0.5pt}\end{center}

\subsection{Esercizio 4}\label{esercizio-4-6}

Trova massimo.

\textbf{Soluzione}

\begin{verbatim}
int max = a[0];
for (int i = 1; i < a.length; i++) {
    if (a[i] > max) {
        max = a[i];
    }
}
\end{verbatim}

\begin{center}\rule{0.5\linewidth}{0.5pt}\end{center}

\subsection{Esercizio 6}\label{esercizio-6-3}

Bubble sort semplice.

\textbf{Soluzione}

\begin{verbatim}
for (int i = 0; i < a.length - 1; i++) {
    for (int j = 0; j < a.length - 1 - i; j++) {
        if (a[j] > a[j+1]) {
            int temp = a[j];
            a[j] = a[j+1];
            a[j+1] = temp;
        }
    }
}
\end{verbatim}

\begin{center}\rule{0.5\linewidth}{0.5pt}\end{center}

\section{9. Metodi}\label{metodi}

\subsection{Esercizio 1}\label{esercizio-1-8}

Metodo quadrato.

\textbf{Soluzione}

\begin{verbatim}
static int quadrato(int x) {
    return x * x;
}
\end{verbatim}

\begin{center}\rule{0.5\linewidth}{0.5pt}\end{center}

\subsection{Esercizio 5}\label{esercizio-5-1}

Metodo che restituisce massimo in un array.

\textbf{Soluzione}

\begin{verbatim}
static int massimo(int[] a) {
    int max = a[0];
    for (int i = 1; i < a.length; i++) {
        if (a[i] > max) {
            max = a[i];
        }
    }
    return max;
}
\end{verbatim}

\begin{center}\rule{0.5\linewidth}{0.5pt}\end{center}

\subsection{Esercizio 6}\label{esercizio-6-4}

Metodo che verifica se un numero è primo.

\textbf{Soluzione}

\begin{verbatim}
static boolean primo(int n) {
    if (n < 2) return false;
    for (int i = 2; i <= Math.sqrt(n); i++) {
        if (n % i == 0) return false;
    }
    return true;
}
\end{verbatim}

\begin{center}\rule{0.5\linewidth}{0.5pt}\end{center}

\section{10. Classi e oggetti}\label{classi-e-oggetti}

\subsection{Esercizio 1}\label{esercizio-1-9}

Classe Studente.

\textbf{Soluzione}

\begin{verbatim}
public class Studente {
    String nome;
    int eta;
}
\end{verbatim}

\begin{center}\rule{0.5\linewidth}{0.5pt}\end{center}

\subsection{Esercizio 4}\label{esercizio-4-7}

Costruttore e instanziazione.

\textbf{Soluzione}

\begin{verbatim}
public class Studente {
    String nome;
    int eta;

    public Studente(String n, int e) {
        nome = n;
        eta = e;
    }
}
\end{verbatim}

\begin{center}\rule{0.5\linewidth}{0.5pt}\end{center}

\subsection{Esercizio 7}\label{esercizio-7}

Classe Prodotto con calcolo totale.

\textbf{Soluzione}

\begin{verbatim}
public class Prodotto {
    String nome;
    double prezzo;
    int quantita;

    public Prodotto(String n, double p, int q) {
        nome = n;
        prezzo = p;
        quantita = q;
    }

    public double totale() {
        return prezzo * quantita;
    }
}
\end{verbatim}

\begin{center}\rule{0.5\linewidth}{0.5pt}\end{center}

\subsection{Esercizio 8}\label{esercizio-8}

Classe Rettangolo.

\textbf{Soluzione}

\begin{verbatim}
public class Rettangolo {
    double base;
    double altezza;

    public Rettangolo(double b, double h) {
        base = b;
        altezza = h;
    }

    public double area() {
        return base * altezza;
    }

    public double perimetro() {
        return 2 * (base + altezza);
    }
}
\end{verbatim}

\begin{center}\rule{0.5\linewidth}{0.5pt}\end{center}

\end{document}
